\documentclass[10.5pt,a4paper]{article}

%% ────────────────────────────────────────────────────────────────────────────
%%  ENCODING & LANGUAGE
%% ────────────────────────────────────────────────────────────────────────────
\usepackage[utf8]{inputenc}
\usepackage[T1]{fontenc}
\usepackage[spanish,es-nodecimaldot]{babel}

%% ────────────────────────────────────────────────────────────────────────────
%%  FUENTE: Helvetica / sans-serif limpia
%% ────────────────────────────────────────────────────────────────────────────
\usepackage{helvet}
\usepackage{lmodern}
\usepackage{microtype}
\renewcommand{\familydefault}{\sfdefault}

%% ────────────────────────────────────────────────────────────────────────────
%%  LAYOUT
%% ────────────────────────────────────────────────────────────────────────────
\usepackage[left=2.6cm, right=2.6cm, top=2.6cm, bottom=3cm]{geometry}
\usepackage{setspace}
\usepackage{parskip}
\setstretch{1.18}
\setlength{\parindent}{0pt}
\setlength{\parskip}{5pt}

%% ────────────────────────────────────────────────────────────────────────────
%%  PAQUETES VISUALES
%% ────────────────────────────────────────────────────────────────────────────
\usepackage{xcolor}
\usepackage{tikz}
\usetikzlibrary{positioning,arrows.meta,shapes.geometric,calc,backgrounds}
\usepackage[most]{tcolorbox}
\usepackage{graphicx}
\usepackage{caption}

%% ────────────────────────────────────────────────────────────────────────────
%%  TABLAS
%% ────────────────────────────────────────────────────────────────────────────
\usepackage{booktabs}
\usepackage{tabularx}
\usepackage{array}
\usepackage{longtable}
\usepackage{colortbl}
\renewcommand{\arraystretch}{1.28}
\setlength{\tabcolsep}{8pt}

%% ────────────────────────────────────────────────────────────────────────────
%%  LISTAS
%% ────────────────────────────────────────────────────────────────────────────
\usepackage{enumitem}

%% ────────────────────────────────────────────────────────────────────────────
%%  SECCIONES Y ENCABEZADOS
%% ────────────────────────────────────────────────────────────────────────────
\usepackage{titlesec}
\usepackage{fancyhdr}
\usepackage{lastpage}

%% ────────────────────────────────────────────────────────────────────────────
%%  HIPERVÍNCULOS
%% ────────────────────────────────────────────────────────────────────────────
\usepackage[hypertexnames=false, hidelinks]{hyperref}
\usepackage{bookmark}


%% ═══════════════════════════════════════════════════════════════════════════
%%  SISTEMA DE COLORES  (design token layer)
%% ═══════════════════════════════════════════════════════════════════════════
%%  Brand
\definecolor{unsaMaroon} {HTML}{7A0018}
\definecolor{unsaMaroonL}{HTML}{F9EFF1}   % maroon 5%
\definecolor{unsaBlue}   {HTML}{1D4ED8}
\definecolor{unsaBlueL}  {HTML}{EFF6FF}   % blue 5%
\definecolor{unsaGold}   {HTML}{D97706}
\definecolor{unsaGoldL}  {HTML}{FFFBEB}   % gold 5%
\definecolor{unsaGreen}  {HTML}{059669}
\definecolor{unsaGreenL} {HTML}{ECFDF5}
%%  Neutral scale
\definecolor{unsaInk}    {HTML}{1E293B}   % headings
\definecolor{unsaBody}   {HTML}{374151}   % body text
\definecolor{unsaSlate}  {HTML}{64748B}   % secondary
\definecolor{unsaMuted}  {HTML}{94A3B8}   % captions
\definecolor{unsaBorder} {HTML}{E2E8F0}   % dividers
\definecolor{unsaLight}  {HTML}{F8FAFC}   % surface
\definecolor{unsaWhite}  {HTML}{FFFFFF}
%%  Cover dark
\definecolor{coverBg}    {HTML}{0F172A}   % slate-900
\definecolor{coverBg2}   {HTML}{1E293B}   % slate-800

%% ─ Configuración de caption y listas (después de definir colores) ──────────
\captionsetup{font=small, labelfont=bf, labelsep=period, skip=6pt}
\setlist[itemize]{leftmargin=1.2cm, itemsep=3pt, topsep=4pt,
  label={\color{unsaMaroon}\rule[0.2ex]{5pt}{5pt}}}
\setlist[enumerate]{leftmargin=1.2cm, itemsep=3pt, topsep=4pt}


%% ═══════════════════════════════════════════════════════════════════════════
%%  HYPERSETUP
%% ═══════════════════════════════════════════════════════════════════════════
\hypersetup{
  colorlinks=true,
  linkcolor=unsaMaroon,
  urlcolor=unsaBlue,
  citecolor=unsaMaroon,
  pdftitle={UNSAP — Whitepaper de Inversión 2026},
  pdfauthor={Equipo UNSAP},
  pdfsubject={EdTech SuperApp universitaria — estrategia de inversión},
  pdfkeywords={UNSAP, EdTech, Flutter, Supabase, Whitepaper, Inversión}
}


%% ═══════════════════════════════════════════════════════════════════════════
%%  FORMATO DE SECCIONES
%% ═══════════════════════════════════════════════════════════════════════════

%% Section — badge numerado + línea inferior
\titleformat{\section}
  {\Large\bfseries\color{unsaInk}}
  {\colorbox{unsaMaroon}{\color{white}\hspace{4pt}\thesection\hspace{4pt}}}
  {0.75em}{}
  [{\color{unsaBorder}\titlerule[0.5pt]}]
\titlespacing*{\section}{0pt}{3ex}{1.5ex}

%% Subsection — número en maroon, texto en ink
\titleformat{\subsection}
  {\normalsize\bfseries\color{unsaInk}}
  {\color{unsaMaroon}\thesubsection.}
  {0.5em}{}[]
\titlespacing*{\subsection}{0pt}{2ex}{0.8ex}

%% Subsubsection
\titleformat{\subsubsection}
  {\small\bfseries\color{unsaSlate}}
  {}
  {0em}{}[]
\titlespacing*{\subsubsection}{0pt}{1.4ex}{0.4ex}


%% ═══════════════════════════════════════════════════════════════════════════
%%  CABECERA / PIE DE PÁGINA
%% ═══════════════════════════════════════════════════════════════════════════
\pagestyle{fancy}
\fancyhf{}
\setlength{\headheight}{16pt}

\fancyhead[L]{%
  {\footnotesize\bfseries\color{unsaMaroon}UNSAP}%
  \enspace{\color{unsaBorder}\textbar}\enspace
  {\footnotesize\color{unsaSlate}Whitepaper de Inversión 2026}%
}
\fancyhead[R]{%
  {\footnotesize\color{unsaSlate}\thepage\ \textsl{de}\ \pageref{LastPage}}%
}
\renewcommand{\headrulewidth}{0.4pt}
\renewcommand{\headrule}{%
  \hbox to\headwidth{\color{unsaBorder}\leaders\hrule height\headrulewidth\hfill}%
}
\renewcommand{\footrulewidth}{0pt}


%% ═══════════════════════════════════════════════════════════════════════════
%%  COMPONENTES TCOLORBOX
%% ═══════════════════════════════════════════════════════════════════════════

%% ─ keybox: alerta / nota clave con borde izquierdo maroon ──────────────────
\newtcolorbox{keybox}{
  enhanced, breakable,
  colback=unsaMaroonL,
  colframe=unsaMaroon,
  boxrule=0pt,
  borderline west={3pt}{0pt}{unsaMaroon},
  arc=0pt, sharp corners,
  left=10pt, right=10pt, top=7pt, bottom=7pt
}

%% ─ signalbox: hipótesis / insight con borde blue ───────────────────────────
\newtcolorbox{signalbox}{
  enhanced, breakable,
  colback=unsaBlueL,
  colframe=unsaBlue,
  boxrule=0pt,
  borderline west={3pt}{0pt}{unsaBlue},
  arc=0pt, sharp corners,
  left=10pt, right=10pt, top=7pt, bottom=7pt
}

%% ─ successbox: hito positivo ───────────────────────────────────────────────
\newtcolorbox{successbox}{
  enhanced, breakable,
  colback=unsaGreenL,
  colframe=unsaGreen,
  boxrule=0pt,
  borderline west={3pt}{0pt}{unsaGreen},
  arc=0pt, sharp corners,
  left=10pt, right=10pt, top=7pt, bottom=7pt
}

%% ─ darkbox: caja oscura para datos de inversión ────────────────────────────
\newtcolorbox{darkbox}{
  colback=coverBg2,
  colframe=coverBg2,
  boxrule=0pt,
  arc=4pt,
  left=14pt, right=14pt, top=12pt, bottom=12pt,
  fontupper=\color{white}
}


%% ═══════════════════════════════════════════════════════════════════════════
%%  COMANDO TABLA — header coloreado
%% ═══════════════════════════════════════════════════════════════════════════
\newcommand{\tableheadrow}{\rowcolor{unsaMaroon}}
\newcommand{\tablehead}[1]{\color{white}\bfseries #1}


%% ═══════════════════════════════════════════════════════════════════════════
%%  DOCUMENTO
%% ═══════════════════════════════════════════════════════════════════════════
\begin{document}

%% ┌──────────────────────────────────────────────────────────────────────────┐
%% │  PORTADA                                                                 │
%% └──────────────────────────────────────────────────────────────────────────┘
\begin{titlepage}
\thispagestyle{empty}

%% Fondo oscuro completo + decoración geométrica
\begin{tikzpicture}[remember picture, overlay]
  %% Background
  \fill[coverBg] (current page.north west) rectangle (current page.south east);
  %% Círculos decorativos sutiles
  \fill[unsaMaroon, opacity=0.18]
    ([xshift=15cm, yshift=-6cm]current page.north west) circle (9cm);
  \fill[unsaBlue, opacity=0.08]
    ([xshift=-2cm, yshift=-20cm]current page.north west) circle (8cm);
  \fill[white, opacity=0.02]
    ([xshift=18cm, yshift=-24cm]current page.north west) circle (10cm);
  %% Barra lateral izquierda maroon
  \fill[unsaMaroon]
    (current page.north west)
    rectangle ([xshift=0.45cm]current page.south west);
  %% Línea dorada separadora (inferior)
  \fill[unsaGold]
    ([yshift=5.8cm]current page.south west)
    rectangle ([yshift=6.1cm]current page.south east);
  %% Banda inferior oscura para metadatos
  \fill[coverBg2]
    (current page.south west)
    rectangle ([yshift=5.8cm]current page.south east);
\end{tikzpicture}

%% ── Institución (zona superior) ────────────────────────────────────────────
\vspace*{1.4cm}
\hspace{0.4cm}
{\color{unsaMuted}\footnotesize\bfseries\uppercase{Universidad Nacional de San Agustín}}
\hspace{0.25cm}
{\color{unsaMaroon!60}\textbar}
\hspace{0.25cm}
{\color{unsaMuted}\footnotesize Arequipa, Perú · 2026}

\vspace{3.5cm}

%% ── Logotipo principal ─────────────────────────────────────────────────────
\begin{center}
  {\fontsize{72}{80}\selectfont\bfseries\color{white}UNSAP}\\[0.35cm]
  {\color{unsaGold}\rule{8cm}{1.2pt}}\\[0.5cm]
  {\large\color{white!80!coverBg}Sistema Operativo de la Vida Universitaria}\\[0.25cm]
  {\normalsize\color{unsaMuted}SuperApp EdTech · Flutter · Supabase · AI-Powered}
\end{center}

\vspace{1.4cm}

%% ── Badge tipo de documento ────────────────────────────────────────────────
\begin{center}
  \begin{tikzpicture}
    \node[fill=unsaMaroon, text=white, inner xsep=14pt, inner ysep=6pt,
          rounded corners=3pt, font=\small\bfseries] {
      WHITEPAPER DE INVERSIÓN
      \enspace\textbar\enspace
      RONDA PRE-SEED / SEED
    };
  \end{tikzpicture}
\end{center}

\vspace{2cm}

%% ── KPI Cards (3 métricas clave) ───────────────────────────────────────────
\begin{center}
\begin{tikzpicture}
  %% Card 1
  \node[fill=white, opacity=0.06, rounded corners=4pt,
        minimum width=4.6cm, minimum height=2.5cm] at (-5.2,0) {};
  \node[text=unsaMuted, font=\footnotesize\bfseries] at (-5.2, 0.6)
        {TAM GLOBAL};
  \node[text=white, font=\fontsize{26}{28}\selectfont\bfseries] at (-5.2, 0)
        {32M};
  \node[text=unsaMuted, font=\scriptsize] at (-5.2,-0.65)
        {estudiantes universitarios · LATAM};
  \draw[unsaMaroon!60, line width=0.5pt, rounded corners=4pt]
        (-7.5,-1.3) rectangle (-2.9,1.3);

  %% Card 2
  \node[fill=unsaMaroon, opacity=0.25, rounded corners=4pt,
        minimum width=4.6cm, minimum height=2.5cm] at (0,0) {};
  \node[text=unsaMuted, font=\footnotesize\bfseries] at (0, 0.6)
        {SOM INICIAL};
  \node[text=white, font=\fontsize{26}{28}\selectfont\bfseries] at (0, 0)
        {70K};
  \node[text=unsaMuted, font=\scriptsize] at (0,-0.65)
        {usuarios · UNSA + Arequipa};
  \draw[unsaMaroon, line width=1pt, rounded corners=4pt]
        (-2.3,-1.3) rectangle (2.3,1.3);

  %% Card 3
  \node[fill=white, opacity=0.06, rounded corners=4pt,
        minimum width=4.6cm, minimum height=2.5cm] at (5.2,0) {};
  \node[text=unsaMuted, font=\footnotesize\bfseries] at (5.2, 0.6)
        {RONDA SEED};
  \node[text=white, font=\fontsize{26}{28}\selectfont\bfseries] at (5.2, 0)
        {\$200K};
  \node[text=unsaMuted, font=\scriptsize] at (5.2,-0.65)
        {USD · runway 18 meses};
  \draw[unsaBlue!60, line width=0.5pt, rounded corners=4pt]
        (2.9,-1.3) rectangle (7.5,1.3);
\end{tikzpicture}
\end{center}

%% ── Metadatos (zona inferior) ──────────────────────────────────────────────
\vspace*{\fill}

\begin{minipage}[b]{0.5\textwidth}
  {\color{unsaMuted}\scriptsize
    \textbf{\color{white!70!coverBg}Versión del producto} \hfill 1.4.0\\[3pt]
    \textbf{\color{white!70!coverBg}Preparado por} \hfill CEO \& Founder\\[3pt]
    \textbf{\color{white!70!coverBg}Stack} \hfill Flutter · Supabase · Dart Isolates\\[3pt]
    \textbf{\color{white!70!coverBg}Fecha} \hfill Marzo 2026
  }
\end{minipage}
\hfill
\begin{minipage}[b]{0.44\textwidth}
  \begin{tikzpicture}
    \node[fill=unsaMaroon!80, text=white, inner xsep=10pt, inner ysep=7pt,
          rounded corners=3pt, font=\scriptsize, align=center] {
      \textbf{DOCUMENTO CONFIDENCIAL}\\[2pt]
      Uso exclusivo para evaluación\\de inversión y due diligence.
    };
  \end{tikzpicture}
\end{minipage}

\vspace{0.8cm}

\end{titlepage}


%% ┌──────────────────────────────────────────────────────────────────────────┐
%% │  ÍNDICES                                                                 │
%% └──────────────────────────────────────────────────────────────────────────┘
\pagenumbering{roman}

{
  \hypersetup{linkcolor=unsaInk}
  \tableofcontents
}
\newpage

\listoftables
\listoffigures
\newpage

\pagenumbering{arabic}


%% ════════════════════════════════════════════════════════════════════════════
\section{Resumen Ejecutivo}
%% ════════════════════════════════════════════════════════════════════════════

La educación superior en Latinoamérica presenta una brecha visible entre la calidad de las experiencias digitales de consumo masivo y la experiencia de uso de portales académicos institucionales. UNSAP responde a esta brecha con una \textbf{SuperApp universitaria} que integra información académica, analítica de rendimiento, red social contextual y gamificación.

El valor diferencial del producto radica en cuatro pilares técnicos y de producto:

\begin{itemize}
  \item Experiencia \textbf{offline-first} con consulta inmediata de información académica.
  \item Arquitectura \textbf{agnóstica al backend institucional} mediante scraping estructurado.
  \item Módulo social y gamificado para \textbf{retención recurrente} del usuario.
  \item Camino claro de monetización \textbf{B2C, B2B2C y B2B} desde el primer año.
\end{itemize}

\begin{signalbox}
\textbf{Hipótesis de inversión:} consolidar UNSAP en el mercado inicial (UNSA), acelerar módulos de IA y marketplace, y empaquetar la arquitectura para expansión regional y licenciamiento white-label a universidades de la región.
\end{signalbox}

\vspace{0.4cm}

%% ── Strip de métricas del resumen ───────────────────────────────────────────
\begin{darkbox}
\begin{tabular}{@{} p{3.5cm} p{3.5cm} p{3.5cm} p{3.5cm} @{}}
  {\color{unsaGold}\footnotesize\bfseries AÑO 1 MAU}  &
  {\color{unsaGold}\footnotesize\bfseries INGRESOS AÑO 3} &
  {\color{unsaGold}\footnotesize\bfseries EBITDA AÑO 3} &
  {\color{unsaGold}\footnotesize\bfseries MODELOS} \\[4pt]
  {\large\bfseries 25,000}    &
  {\large\bfseries \$1.02M}   &
  {\large\bfseries \$624K}    &
  {\large\bfseries 3} \\[2pt]
  {\scriptsize\color{unsaMuted}usuarios activos} &
  {\scriptsize\color{unsaMuted}ingresos brutos USD} &
  {\scriptsize\color{unsaMuted}referencial USD} &
  {\scriptsize\color{unsaMuted}B2C · B2B2C · B2B}
\end{tabular}
\end{darkbox}


%% ════════════════════════════════════════════════════════════════════════════
\section{Problema Estructural y Oportunidad}
%% ════════════════════════════════════════════════════════════════════════════

\subsection{Fricción digital actual}

Los sistemas académicos tradicionales presentan latencias altas en picos de uso, interfaces de baja usabilidad y dependencia fuerte de conectividad. Esto genera ansiedad operativa y bajo compromiso de uso entre los estudiantes.

\subsection{Fragmentación del ciclo de vida estudiantil}

La información académica, la colaboración entre pares y la vida social estudiantil se distribuyen en plataformas no integradas. La consecuencia es pérdida de contexto, baja productividad y escasa trazabilidad de decisiones académicas.

\subsection{Carga cognitiva y ausencia de predicción}

Sin herramientas predictivas integradas, el estudiante realiza cálculos manuales para tomar decisiones críticas sobre su rendimiento. Esta fricción representa una oportunidad para soluciones con analítica accionable en tiempo real.


%% ════════════════════════════════════════════════════════════════════════════
\section{Solución UNSAP y Ventaja Competitiva}
%% ════════════════════════════════════════════════════════════════════════════

\subsection{Arquitectura offline-first y experiencia de cero fricción}

UNSAP persiste información en almacenamiento local (Hive), permitiendo consulta instantánea de notas, horario y estadísticas, incluso sin red o ante degradación de servicios externos.

\subsection{Motor de extracción agnóstico}

El sistema utiliza un motor de scraping en isolates para capturar y estructurar datos desde portales legacy sin bloquear la interfaz de usuario.

\subsection{Simulador académico O(1)}

Se incorporan cálculos de proyección de notas y ponderados con operaciones optimizadas para respuestas inmediatas, mejorando la toma de decisiones del estudiante.

\begin{figure}[htbp]
\centering
\resizebox{0.98\textwidth}{!}{%
\begin{tikzpicture}[
  node distance=14mm,
  source/.style={rectangle, rounded corners=3mm, draw=unsaMaroon,
    fill=unsaMaroonL, align=center, minimum height=1.5cm,
    minimum width=3.8cm, line width=0.7pt, font=\small},
  process/.style={rectangle, rounded corners=3mm, draw=unsaBlue,
    fill=unsaBlueL, align=center, minimum height=1.5cm,
    minimum width=4.8cm, line width=0.7pt, font=\small},
  output/.style={rectangle, rounded corners=3mm, draw=unsaGreen,
    fill=unsaGreenL, align=center, minimum height=1.4cm,
    minimum width=3.6cm, line width=0.7pt, font=\small},
  arr/.style={-{Latex[length=2.8mm, width=2mm]}, draw=unsaSlate, line width=1pt}
]
  \node[source] (legacy)  {\textbf{Portal Universitario}\\\textit{HTML legacy}};
  \node[process, right=16mm of legacy] (scraper)
        {\textbf{Isolate Scraper Engine}\\parsing background\\normalización JSON};
  \node[source, right=16mm of scraper] (hive)
        {\textbf{Capa UNSAP}\\Hive Local Storage};

  \node[output, below left=14mm and -5mm of hive]   (sim)   {\textbf{Simulador}\\\textit{Proyección notas}};
  \node[output, below=14mm of hive]                  (chart) {\textbf{Analítica}\\\textit{fl\_Chart}};
  \node[output, below right=14mm and -5mm of hive]   (panel) {\textbf{Panel de Notas}\\\textit{Dashboard}};

  \draw[arr] (legacy)  -- (scraper);
  \draw[arr] (scraper) -- (hive);
  \draw[arr] (hive) -- (sim);
  \draw[arr] (hive) -- (chart);
  \draw[arr] (hive) -- (panel);
\end{tikzpicture}
}
\caption{Flujo de arquitectura de datos}
\label{fig:flujo-datos}
\end{figure}


%% ════════════════════════════════════════════════════════════════════════════
\section{Motor de Retención: Comunidad y Gamificación}
%% ════════════════════════════════════════════════════════════════════════════

\subsection{Campus social institucional}

UNSAP integra un feed en tiempo real sobre Supabase para publicaciones estudiantiles, interacciones y notificaciones. El anonimato controlado incrementa participación sin perder trazabilidad de calidad de comunidad.

\subsection{Identidad verificada y reputación}

El uso de correo institucional permite elevar niveles de confianza del ecosistema y habilita segmentación de alto valor para futuras verticales de monetización.

\subsection{Medallas local-first y actualizaciones optimistas}

El sistema de medallas con respuesta visual instantánea (0\,ms percibidos en UI) mejora la motivación y recurrencia de uso, mientras la sincronización ocurre en segundo plano con manejo de rollback.

\begin{figure}[htbp]
\centering
\resizebox{0.8\textwidth}{!}{%
\begin{tikzpicture}[
  every node/.style={font=\small},
  ring/.style={circle, draw=unsaMaroon, line width=0.8pt,
               minimum width=3.4cm, align=center,
               fill=unsaLight, inner sep=4pt},
  arr/.style={-{Latex[length=2.8mm, width=2mm]}, draw=unsaSlate, line width=1pt, bend left=20}
]
  \node[ring, fill=unsaMaroonL] (trigger) at (0,3.2)
        {\textbf{Trigger}\\Push · nueva nota\\actividad social};
  \node[ring, fill=unsaBlueL]   (accion)  at (3.6,0)
        {\textbf{Acción}\\Apertura de app\\carga instantánea};
  \node[ring, fill=unsaGoldL]   (reward)  at (0,-3.2)
        {\textbf{Recompensa}\\insights · contenido\\medallas};
  \node[ring, fill=unsaGreenL]  (invest)  at (-3.6,0)
        {\textbf{Inversión}\\perfil · vínculos\\publicación};

  \draw[arr] (trigger) to (accion);
  \draw[arr] (accion)  to (reward);
  \draw[arr] (reward)  to (invest);
  \draw[arr] (invest)  to (trigger);
\end{tikzpicture}
}
\caption{Bucle de retención basado en comportamiento (Hook Model)}
\label{fig:retencion-hook}
\end{figure}


%% ════════════════════════════════════════════════════════════════════════════
\section{Mercado, Posicionamiento y Competencia}
%% ════════════════════════════════════════════════════════════════════════════

\subsection{Dimensionamiento TAM, SAM y SOM}

\begin{table}[htbp]
\centering
\caption{Estimación de mercado objetivo (referencial)}
\begin{tabularx}{\textwidth}{>{\bfseries}p{2.2cm} X X}
\tableheadrow
  \tablehead{Segmento} & \tablehead{Alcance} & \tablehead{Potencial de valor} \\
\midrule
\rowcolor{unsaMaroonL}
TAM & 32 millones de estudiantes universitarios en Latinoamérica &
  Mercado amplio de suscripción, servicios y publicidad educativa \\
\rowcolor{white}
SAM & $\approx$\,1.5 millones en Perú &
  Escala nacional para crecimiento progresivo y alianzas \\
\rowcolor{unsaLight}
SOM & $\approx$\,70,000 en UNSA y Arequipa &
  Validación temprana, iteración rápida y evidencia de PMF \\
\bottomrule
\end{tabularx}
\end{table}

\subsection{Tesis competitiva}

Las plataformas tradicionales (LMS/ERP institucionales) operan bajo lógica B2B. UNSAP adopta un enfoque \textbf{bottom-up}, centrado en el estudiante, con mayor velocidad de iteración y diferenciación en comunidad + gamificación.


%% ════════════════════════════════════════════════════════════════════════════
\section{Modelo de Negocio y Monetización}
%% ════════════════════════════════════════════════════════════════════════════

\begin{table}[htbp]
\centering
\caption{Arquitectura de ingresos propuesta}
\begin{tabularx}{\textwidth}{>{\bfseries}p{2.2cm} X X}
\tableheadrow
  \tablehead{Canal} & \tablehead{Mecánica de monetización} & \tablehead{Valor estratégico} \\
\midrule
\rowcolor{white}
B2C &
  Suscripción \textbf{UNSAP PRO} con funciones avanzadas de analítica, personalización e IA &
  Ingreso recurrente y relación directa con el usuario final \\
\rowcolor{unsaLight}
B2B2C &
  Publicidad nativa y lead generation de alto contexto académico &
  Margen superior por segmentación y precisión de audiencia \\
\rowcolor{white}
B2B &
  Licenciamiento white-label para universidades (setup + fee anual por estudiante) &
  Escalamiento institucional con contratos recurrentes \\
\bottomrule
\end{tabularx}
\end{table}


%% ════════════════════════════════════════════════════════════════════════════
\section{Roadmap Tecnológico 2026--2028}
%% ════════════════════════════════════════════════════════════════════════════

\subsection{Fase 1 — Dominio local y robustez \normalfont\small(Meses 1–6)}
Enfoque en penetración de mercado inicial, hardening técnico, observabilidad y resiliencia operativa.

\subsection{Fase 2 — Servicios de valor \normalfont\small(Meses 7–18)}
Integración de UNSAP GPT para consultas académicas, despliegue de marketplace C2C y capa de transacciones.

\subsection{Fase 3 — Escala regional y B2B \normalfont\small(Meses 19–36)}
Expansión multiuniversidad y consolidación del producto como plataforma modular para nuevas instituciones.

\begin{figure}[htbp]
\centering
\resizebox{0.98\textwidth}{!}{%
\begin{tikzpicture}[x=0.66cm, y=0.9cm]
  %% Grid de fondo
  \fill[unsaLight] (0,0) rectangle (14,4.5);
  \foreach \x in {0,2,...,14}
    {\draw[unsaBorder, line width=0.4pt] (\x,0) -- (\x,4.5);}
  \draw[unsaBorder, line width=0.5pt] (0,0) rectangle (14,4.5);

  %% Etiquetas de fases
  \node[anchor=east, font=\footnotesize\bfseries, text=unsaMaroon] at (-0.3,3.7) {Fase 1};
  \node[anchor=east, font=\footnotesize\bfseries, text=unsaBlue]   at (-0.3,2.3) {Fase 2};
  \node[anchor=east, font=\footnotesize\bfseries, text=unsaGreen]  at (-0.3,0.9) {Fase 3};

  %% Barra Fase 1
  \fill[unsaMaroon, rounded corners=3pt] (0.12,3.2) rectangle (3.88,4.2);
  \node[text=white, font=\footnotesize\bfseries] at (2,3.7) {Q3 2026 — Q4 2026};

  %% Barra Fase 2
  \fill[unsaBlue, rounded corners=3pt] (3.95,1.8) rectangle (9.65,2.8);
  \node[text=white, font=\footnotesize\bfseries] at (6.8,2.3) {Q1 2027 — Q4 2027};

  %% Barra Fase 3
  \fill[unsaGreen, rounded corners=3pt] (9.72,0.4) rectangle (13.88,1.4);
  \node[text=white, font=\footnotesize\bfseries] at (11.8,0.9) {Q1 2028 — Q4 2028};

  %% Línea de tiempo (eje X)
  \foreach \x/\q in {0/Q3'26, 2/Q1'27, 4/Q3'27, 6/Q1'28, 8/Q3'28, 10/Q1'29, 12/Q3'29, 14/Q1'30}
    {\node[anchor=north, font=\scriptsize, text=unsaSlate] at (\x,-0.15) {\q};}
\end{tikzpicture}
}
\caption{Roadmap por fases de crecimiento (Gantt referencial)}
\label{fig:gantt}
\end{figure}


%% ════════════════════════════════════════════════════════════════════════════
\section{Proyecciones Financieras — Escenario Referencial}
%% ════════════════════════════════════════════════════════════════════════════

\begin{table}[htbp]
\centering
\caption{Proyección financiera a 36 meses (USD)}
\begin{tabularx}{\textwidth}{>{\bfseries}X r r r}
\tableheadrow
  \tablehead{Métrica} & \tablehead{Año 1} & \tablehead{Año 2} & \tablehead{Año 3} \\
\midrule
\rowcolor{white}
MAU estimado            & 25,000   & 120,000   & 450,000 \\
\rowcolor{unsaLight}
Ingresos suscripción PRO& \$18,000  & \$86,400  & \$324,000 \\
\rowcolor{white}
Ingresos Ad-Tech \& Lead Gen & \$15,000 & \$110,000 & \$450,000 \\
\rowcolor{unsaLight}
Ingresos SaaS white-label    & —        & \$40,000  & \$250,000 \\
\midrule
\rowcolor{unsaMaroonL}
\bfseries Ingresos brutos estimados  &
  \bfseries\$33,000 & \bfseries\$236,400 & \bfseries\$1,024,000 \\
\rowcolor{white}
OPEX estimado           & \$80,000  & \$150,000 & \$400,000 \\
\rowcolor{unsaGreenL}
\bfseries EBITDA estimado &
  \bfseries\color{unsaMaroon}(\$47,000) & \bfseries\$86,400 & \bfseries\color{unsaGreen}\$624,000 \\
\bottomrule
\end{tabularx}
\end{table}

\begin{successbox}
El modelo proyecta \textbf{break-even operativo en el Año 2} y un EBITDA positivo de \$624K al cierre del Año 3, con una base de 450K usuarios activos en tres países.
\end{successbox}


%% ════════════════════════════════════════════════════════════════════════════
\section{Requerimiento de Capital y Uso de Fondos}
%% ════════════════════════════════════════════════════════════════════════════

Se propone una ronda pre-seed/seed de referencia en torno a \textbf{\$200,000 USD}, orientada a extender runway operativo y ejecutar roadmap de 18 meses.

\begin{table}[htbp]
\centering
\caption{Asignación objetivo de capital}
\begin{tabularx}{\textwidth}{>{\bfseries}p{4.2cm} c X}
\tableheadrow
  \tablehead{Línea de inversión} & \tablehead{Peso} & \tablehead{Aplicación} \\
\midrule
\rowcolor{white}
Ingeniería y producto   & 50\% &
  Equipo técnico senior para arquitectura, desarrollo y calidad \\
\rowcolor{unsaLight}
Adquisición y growth    & 25\% &
  Embajadores, marketing de performance y expansión de base \\
\rowcolor{white}
Infraestructura técnica & 15\% &
  Escalamiento cloud, observabilidad y servicios de soporte \\
\rowcolor{unsaLight}
Legal y estructuración  & 10\% &
  Gobierno corporativo, propiedad intelectual y cumplimiento \\
\bottomrule
\end{tabularx}
\end{table}


%% ════════════════════════════════════════════════════════════════════════════
\section{Riesgos y Mitigación}
%% ════════════════════════════════════════════════════════════════════════════

\begin{table}[htbp]
\centering
\caption{Matriz ejecutiva de riesgos}
\begin{tabularx}{\textwidth}{>{\bfseries}p{2.4cm} X X}
\tableheadrow
  \tablehead{Tipo} & \tablehead{Riesgo principal} & \tablehead{Mitigación y foso defensivo} \\
\midrule
\rowcolor{white}
Técnico &
  Cambios de estructura en portales fuente afectan scraping &
  Arquitectura de adaptadores, alertas tempranas y estrategia de hot-fix \\
\rowcolor{unsaLight}
Operativo &
  Bloqueo por políticas de red o límites institucionales &
  Rate limiting, ventanas de sincronización y resiliencia offline-first \\
\rowcolor{white}
Competitivo &
  Aparición de app oficial institucional &
  Ventaja en velocidad de iteración, comunidad y gamificación persistente \\
\rowcolor{unsaLight}
Regulatorio &
  Exigencias de protección de datos y cumplimiento &
  Gobierno de datos, políticas de consentimiento y trazabilidad de acceso \\
\bottomrule
\end{tabularx}
\end{table}


%% ════════════════════════════════════════════════════════════════════════════
\section{Conclusiones y Tesis de Inversión}
%% ════════════════════════════════════════════════════════════════════════════

UNSAP combina un problema real de alta frecuencia con una solución técnicamente madura y un modelo de ingresos escalable. La plataforma no solo optimiza la experiencia académica, sino que construye un \textbf{activo digital de relación continua} con el estudiante.

\begin{keybox}
\textbf{Bajo una ejecución disciplinada del roadmap} y una gobernanza adecuada de datos y crecimiento, el proyecto presenta condiciones sólidas para financiamiento de etapa temprana con potencial de \textbf{expansión regional y salida estratégica} de mediano plazo.
\end{keybox}

\vspace{0.5cm}

%% ── Resumen visual de tesis ─────────────────────────────────────────────────
\begin{darkbox}
\begin{center}
  {\color{unsaGold}\bfseries POR QUÉ UNSAP AHORA}\\[0.6cm]
  \begin{tabular}{@{} p{4.5cm} p{4.5cm} p{4.5cm} @{}}
    {\color{unsaGold}\small\bfseries PROBLEMA REAL} &
    {\color{unsaGold}\small\bfseries TECH MADURA} &
    {\color{unsaGold}\small\bfseries MERCADO CLARO} \\[4pt]
    {\scriptsize\color{white!80!coverBg}
      Millones de estudiantes con portales deficientes y sin herramientas de decisión académica.} &
    {\scriptsize\color{white!80!coverBg}
      Flutter, Supabase y arquitectura offline-first probada en producción.} &
    {\scriptsize\color{white!80!coverBg}
      TAM de 32M en LATAM con punto de entrada validado en UNSA (70K usuarios).}
  \end{tabular}
\end{center}
\end{darkbox}


%% ════════════════════════════════════════════════════════════════════════════
\appendix
\section{Stack Tecnológico Consolidado}
%% ════════════════════════════════════════════════════════════════════════════

\begin{longtable}{>{\bfseries\color{unsaMaroon}}p{4cm} p{10cm}}
\tableheadrow
  \tablehead{Componente} & \tablehead{Tecnologías} \\
\midrule
\endhead
Frontend móvil          & Flutter 3.9.2, Dart 3.9.2 \\
\rowcolor{unsaLight}
Arquitectura de app     & MVVM, Provider, Clean Architecture \\
Persistencia local      & Hive, flutter\_secure\_storage \\
\rowcolor{unsaLight}
Backend y tiempo real   & Supabase, PostgreSQL, canales Realtime \\
Analítica visual        & fl\_chart y módulos de simulación académica \\
\rowcolor{unsaLight}
Red y parsing           & http, parser HTML, isolates para procesos asíncronos \\
Caché y media           & cached\_network\_image, share\_plus \\
\rowcolor{unsaLight}
Configuración           & flutter\_dotenv y gestión de entorno \\
\bottomrule
\end{longtable}


%% ════════════════════════════════════════════════════════════════════════════
\section{Notas de Alcance y Uso}
%% ════════════════════════════════════════════════════════════════════════════

Las cifras de mercado y proyecciones de este documento se presentan como escenario referencial para modelado y decisión de inversión. Se recomienda validación continua con datos operativos reales a medida que avance la implementación del roadmap.

\vspace{1cm}

\begin{center}
  {\color{unsaBorder}\rule{0.6\textwidth}{0.5pt}}\\[0.4cm]
  {\scriptsize\color{unsaMuted}
    \textbf{UNSAP} — Documento Confidencial de Inversión
    \enspace·\enspace
    Versión 1.4.0
    \enspace·\enspace
    Marzo 2026
    \enspace·\enspace
    Arequipa, Perú
  }
\end{center}

\end{document}
